\section{Síntesis y ampliación}

\subsection{Repaso de los puntos clave}

En esta clase se derivó detalladamente el algoritmo de entrenamiento y muestreo de DDPM; su contenido central puede resumirse en lo siguiente:

\begin{enumerate}
    \item El \textbf{proceso hacia adelante} es un proceso predefinido de inyección de ruido gaussiano, con una distribución de salto de forma cerrada $q(x_t|x_0)$, que permite el muestreo en un solo paso.
    \item La \textbf{descomposición del ELBO} transforma el objetivo de entrenamiento para que coincida con la distribución posterior real de transición en cada paso temporal.
    \item La \textbf{posterior real} $q(x_{t-1}|x_t, x_0)$ es una distribución gaussiana cuya varianza puede calcularse directamente y cuyo valor medio depende de $x_0$.
    \item Tres \textbf{parametrizaciones}: predicción del valor medio, predicción de $x_0$ y predicción de ruido son matemáticamente equivalentes; en la práctica, la \textbf{predicción de ruido} es la más estable.
    \item \textbf{Objetivo simplificado de entrenamiento}: $L_{\text{simple}} = \mathbb{E}\|\epsilon - \epsilon_\theta(x_t, t)\|^2$, elegante y eficiente.
    \item La \textbf{capacidad de expresión} proviene de la no linealidad introducida por la red neuronal en el proceso inverso: gaussiano paso a paso pero globalmente no gaussiano.
    \item Existe una conexión profunda entre la \textbf{predicción de ruido} y los modelos basados en puntuación (score-based), sentando las bases para un marco teórico unificado.
\end{enumerate}

\subsection{Anticipación de las próximas clases}

El ponente mencionará que se discutirán los siguientes temas:
\begin{itemize}
    \item La conexión formal entre predicción de ruido y score matching.
    \item El marco Score SDE: extensión del tiempo discreto de difusión al continuo.
    \item Diferentes definiciones del proceso hacia adelante.
    \item Métodos de muestreo más rápidos (como DDIM).
\end{itemize}

\subsection{Lecturas adicionales}

\begin{itemize}
    \item \textbf{Artículo original de DDPM}: Ho, Jain, Abbeel. ``Denoising Diffusion Probabilistic Models.'' NeurIPS 2020.
    \item \textbf{DDPM mejorado}: Nichol, Dhariwal. ``Improved Denoising Diffusion Probabilistic Models.'' ICML 2021. Explora estrategias para aprender la varianza.
    \item \textbf{Score SDE}: Song et al. ``Score-Based Generative Modeling through Stochastic Differential Equations.'' ICLR 2021. Unifica los modelos DDPM y score-based.
    \item \textbf{DDIM}: Song, Meng, Ermon. ``Denoising Diffusion Implicit Models.'' ICLR 2021. Propone un método de muestreo determinista que reduce drásticamente el número de pasos de muestreo.
    \item \textbf{P2 Weighting}: Choi et al. ``Perception Prioritized Training of Diffusion Models.'' CVPR 2022. Explora estrategias adaptativas de ponderación del paso temporal.
    \item \textbf{v-prediction}: Salimans, Ho. ``Progressive Distillation for Fast Sampling of Diffusion Models.'' ICLR 2022. Propone la parametrización v-prediction, que combina las ventajas de la predicción de $x_0$ y la predicción de $\epsilon$.
\end{itemize}

\end{document}
